% ===========================================================================
\section{Supplementary Note S4: the benchmark's implied buffalo factor}
\label{app:buffalo}
% ===========================================================================

The South Asian ammonia inventory we benchmark against \citep{sanh2026} reports implicit emission
factors by livestock type, but its table omits buffalo --- the species its own text names as the
reason Pakistan's grazing emissions are high, and the species that dominates Pakistani nitrogen
excretion. Its released dataset is gridded emissions only \citep{sanh2026data}, so the underlying
animal masses, excretion rates and management shares cannot be recovered from it directly. A residual
back-calculation is therefore the only available route, and this note sets out what it does and does
not establish.

\paragraph{Method.} We apply the benchmark's own tabulated factors to the same FAOSTAT 2015
populations our historical series uses, subtract the result from its published Pakistan livestock
total (grazing plus manure management, \SI{495.7}{}~kt), and attribute the residual to the animals its
table does not cover.

\begin{table}[!htbp]
\centering
\footnotesize
\begin{tabular}{@{}lr@{}}
\toprule
\textbf{Term} & \textbf{kt \NHthree{}} \\
\midrule
Benchmark Pakistan livestock total (grazing + manure management) & 495.7 \\
\quad less broiler chickens & $-$102.4 \\
\quad less dairy cattle & $-$77.0 \\
\quad less goats & $-$75.9 \\
\quad less non-dairy cattle & $-$60.2 \\
\quad less sheep & $-$42.4 \\
\midrule
\textbf{Residual, unattributed by the published table} & \textbf{137.8} \\
Unattributed population (buffalo + camels) & 36.6 M head \\
\midrule
Implied factor over buffalo and camels & 3.76 \\
\textbf{Implied buffalo factor, upper bound} & \textbf{3.87} \\
Their tabulated dairy-cattle factor, for scale & 6.33 \\
\bottomrule
\end{tabular}
\caption{Back-calculated buffalo emission factor (kg~\NHthree{} per head per year). \emph{This is our
derivation, not a published value.} The residual also absorbs camels and any category the published
table omits, so it is an upper bound on the buffalo factor alone. Generated by
\texttt{src/pkinv/eval/buffalo\_implied\_ef.py}.}
\label{tab:buffalo}
\end{table}

\paragraph{What it indicates.} An implied factor of at most \SI{3.9}{\kilogram} per head sits
\emph{below} the benchmark's own dairy-cattle value of \SI{6.33}{}, for an animal with substantially
higher nitrogen excretion. The only configuration that reproduces such a value is one in which
buffalo manure is overwhelmingly deposited at pasture, where ammonia losses are small, rather than
passing through housing, storage and field application. Our own species-resolved chain gives
\SI{9.83}{} per head and the GLEAM~3.0 South Asian management split --- which assumes no buffalo
grazing at all --- gives \SI{11.69}{}. Our value therefore lies between two external
parameterisations that disagree with each other by a factor of three.

\paragraph{What it does not establish.} The calculation is a residual, so it inherits every
difference between our activity data and theirs, not only the management split. It cannot separate
buffalo from camels. And the two external parameterisations that bracket our value are not
statistically independent --- both draw on overlapping FAO activity data --- so they bound the
plausible range rather than corroborating each other. The conclusion this supports is that the
manure-management partition is the most likely locus of the residual livestock divergence, which is
why a national manure-management survey is identified in the main text's limitations as the
highest-value data collection this work points to. It is not an attribution, and no result in the
main text depends on it.
